% GENERATED FILE. Edit canonical lessons, then run npm run generate:books.
% canonical-source-hash: fnv1a64:6ae4f1e956cf3169
% canonical-lessons: GU-C352-sarkar, GU-C352-rajya, GU-C352-satta, GU-C352-praja, GU-C352-ghadvu

\chapter{Who Governs}
\label{ch:gu-who-governs}
\hlchaptermodality{352}

\begin{chapteropening}
\textbf{By the end of this chapter:} \emph{I can say who makes the rules.}

Complete the last lesson of chapter 352: I can say who makes the rules.
\end{chapteropening}

\section[\texorpdfstring{\gu{સરકાર} (sarkār)}{સરકાર (sarkār)}]{\gu{સરકાર} (sarkār) — a government}

\label{lesson:GU-C352-sarkar}



\begin{quote}
Before the new one: say the Gujarati for true to the facts, then the Gujarati for to be taken in.
\end{quote}

\subsection*{You'll want to know: \gu{સરકાર}}
\textbf{\gu{સરકાર}} — \emph{sarkār} — \textquotedblleft{}a government\textquotedblright{}.

\begin{quote}
\gu{સરકારે} \gu{ગામમાં} \gu{નવી} \gu{શાળા} \gu{બનાવી}.
\end{quote}

\begin{quote}
\emph{sarkāre gāmmā̃ navī śāḷā banāvī.}
\end{quote}

\begin{quote}
The government built a new school in the village.
\end{quote}

A feminine noun.

\begin{cousinweb}
\textbf{\gu{સરકાર}} (\emph{sarkār}) is usually traced to Persian \emph{sarkār}, \textquotedblleft{}head of affairs\textquotedblright{}, built from \emph{sar}, \textquotedblleft{}head\textquotedblright{}, and \emph{kār}, \textquotedblleft{}work\textquotedblright{}.
\end{cousinweb}

\begin{grammarlens}[title={What you've built}]
One more word for giving reasons and opinions.
\end{grammarlens}

\subsection*{Your turn}
\begin{itemize}
\raggedright
  \item \emph{Say it:} \emph{sarkār}
  \item \emph{Say it:} \emph{sarkār}, once more
  \item \emph{Say it:} \emph{sarkāre gāmmā̃ navī śāḷā banāvī.} — The government built a new school in the village.
  \item \emph{Recall:} say the Gujarati for true to the facts, then the Gujarati for to be taken in, then say \emph{sarkār} again
\end{itemize}

\begin{tcolorbox}[breakable,colback=teal!4,colframe=teal!35!black,title={Before you move on}]
What does \gu{સરકાર} mean? (\textquotedblleft{}A government\textquotedblright{}.) Say it once more.
\end{tcolorbox}
\section[\texorpdfstring{\gu{રાજ્ય} (rājya)}{રાજ્ય (rājya)}]{\gu{રાજ્ય} (rājya) — a state (of a country)}

\label{lesson:GU-C352-rajya}



\begin{quote}
Before the new one: say the Gujarati for to be taken in, then the Gujarati for a government.
\end{quote}

\subsection*{You'll want to know: \gu{રાજ્ય}}
\textbf{\gu{રાજ્ય}} — \emph{rājya} — \textquotedblleft{}a state (of a country)\textquotedblright{}.

\begin{quote}
\gu{ગુજરાત} \gu{એક} \gu{મોટું} \gu{રાજ્ય} \gu{છે}.
\end{quote}

\begin{quote}
\emph{gujrāt ek moṭũ rājya che.}
\end{quote}

\begin{quote}
Gujarat is a big state.
\end{quote}

A neuter noun.

\begin{cousinweb}
\textbf{\gu{રાજ્ય}} (\emph{rājya}) is a tatsama, borrowed from Sanskrit \emph{rājya}, \textquotedblleft{}kingship, a kingdom\textquotedblright{} (Monier-Williams).
\end{cousinweb}

\begin{grammarlens}[title={What you've built}]
One more word for giving reasons and opinions.
\end{grammarlens}

\subsection*{Your turn}
\begin{itemize}
\raggedright
  \item \emph{Say it:} \emph{rājya}
  \item \emph{Say it:} \emph{rājya}, once more
  \item \emph{Say it:} \emph{gujrāt ek moṭũ rājya che.} — Gujarat is a big state.
  \item \emph{Recall:} say the Gujarati for to be taken in, then the Gujarati for a government, then say \emph{rājya} again
\end{itemize}

\begin{tcolorbox}[breakable,colback=teal!4,colframe=teal!35!black,title={Before you move on}]
What does \gu{રાજ્ય} mean? (\textquotedblleft{}A state (of a country)\textquotedblright{}.) Say it once more.
\end{tcolorbox}
\section[\texorpdfstring{\gu{સત્તા} (sattā)}{સત્તા (sattā)}]{\gu{સત્તા} (sattā) — power to rule}

\label{lesson:GU-C352-satta}



\begin{quote}
Before the new one: say the Gujarati for a government, then the Gujarati for a state (of a country).
\end{quote}

\subsection*{You'll want to know: \gu{સત્તા}}
\textbf{\gu{સત્તા}} — \emph{sattā} — \textquotedblleft{}power to rule\textquotedblright{}.

\begin{quote}
\gu{નિયમ} \gu{બદલવાની} \gu{સત્તા} \gu{ફક્ત} \gu{સરકાર} \gu{પાસે} \gu{છે}.
\end{quote}

\begin{quote}
\emph{niyam badalvānī sattā phakt sarkār pāse che.}
\end{quote}

\begin{quote}
Only the government has the power to change the rules.
\end{quote}

A feminine noun.

\begin{cousinweb}
\textbf{\gu{સત્તા}} (\emph{sattā}) is borrowed from Sanskrit \emph{sattā}, \textquotedblleft{}existence, being\textquotedblright{} (Monier-Williams); \textquotedblleft{}power, authority\textquotedblright{} is the sense it has taken in Gujarati.
\end{cousinweb}

\begin{grammarlens}[title={What you've built}]
One more word for giving reasons and opinions.
\end{grammarlens}

\subsection*{Your turn}
\begin{itemize}
\raggedright
  \item \emph{Say it:} \emph{sattā}
  \item \emph{Say it:} \emph{sattā}, once more
  \item \emph{Say it:} \emph{niyam badalvānī sattā phakt sarkār pāse che.} — Only the government has the power to change the rules.
  \item \emph{Recall:} say the Gujarati for a government, then the Gujarati for a state (of a country), then say \emph{sattā} again
\end{itemize}

\begin{tcolorbox}[breakable,colback=teal!4,colframe=teal!35!black,title={Before you move on}]
What does \gu{સત્તા} mean? (\textquotedblleft{}Power to rule\textquotedblright{}.) Say it once more.
\end{tcolorbox}
\section[\texorpdfstring{\gu{પ્રજા} (prajā)}{પ્રજા (prajā)}]{\gu{પ્રજા} (prajā) — the people of a state}

\label{lesson:GU-C352-praja}



\begin{quote}
Before the new one: say the Gujarati for a state (of a country), then the Gujarati for power to rule.
\end{quote}

\subsection*{You'll want to know: \gu{પ્રજા}}
\textbf{\gu{પ્રજા}} — \emph{prajā} — \textquotedblleft{}the people of a state\textquotedblright{}.

\begin{quote}
\gu{સારી} \gu{સરકાર} \gu{પ્રજાની} \gu{વાત} \gu{સાંભળે} \gu{છે}.
\end{quote}

\begin{quote}
\emph{sārī sarkār prajānī vāt sā̃bhḷe che.}
\end{quote}

\begin{quote}
A good government listens to its people.
\end{quote}

A feminine noun.

\begin{cousinweb}
\textbf{\gu{પ્રજા}} (\emph{prajā}) is a tatsama, borrowed from Sanskrit \emph{prajā}, \textquotedblleft{}offspring; subjects\textquotedblright{} (Monier-Williams).
\end{cousinweb}

\begin{grammarlens}[title={What you've built}]
One more word for giving reasons and opinions.
\end{grammarlens}

\subsection*{Your turn}
\begin{itemize}
\raggedright
  \item \emph{Say it:} \emph{prajā}
  \item \emph{Say it:} \emph{prajā}, once more
  \item \emph{Say it:} \emph{sārī sarkār prajānī vāt sā̃bhḷe che.} — A good government listens to its people.
  \item \emph{Recall:} say the Gujarati for a state (of a country), then the Gujarati for power to rule, then say \emph{prajā} again
\end{itemize}

\begin{tcolorbox}[breakable,colback=teal!4,colframe=teal!35!black,title={Before you move on}]
What does \gu{પ્રજા} mean? (\textquotedblleft{}The people of a state\textquotedblright{}.) Say it once more.
\end{tcolorbox}
\section[\texorpdfstring{\gu{ઘડવું} (ghaḍvũ)}{ઘડવું (ghaḍvũ)}]{\gu{ઘડવું} (ghaḍvũ) — to shape; to frame (a law or a plan)}

\label{lesson:GU-C352-ghadvu}



\begin{quote}
Before the new one: say the Gujarati for power to rule, then the Gujarati for the people of a state.
\end{quote}

\subsection*{You'll want to know: \gu{ઘડવું}}
\textbf{\gu{ઘડવું}} — \emph{ghaḍvũ} — \textquotedblleft{}to shape; to frame (a law or a plan)\textquotedblright{}.

\begin{quote}
\gu{સરકાર} \gu{નવો} \gu{કાયદો} \gu{ઘડશે}.
\end{quote}

\begin{quote}
\emph{sarkār navo kāydo ghaḍśe.}
\end{quote}

\begin{quote}
The government will draw up a new law.
\end{quote}

Past, with \gu{તેણે} (he or she …): \gu{તેણે} \gu{નવો} \gu{નિયમ} \gu{ઘડ્યો}, \emph{teṇe navo niyam ghaḍyo}.

\begin{cousinweb}
\textbf{\gu{ઘડવું}} (\emph{ghaḍvũ}) is inherited from Sanskrit \emph{gháṭatē}; CDIAL glosses the Gujarati verb \textquotedblleft{}to shape, to hammer\textquotedblright{} (CDIAL 4407). A law is shaped as a smith shapes metal.
\end{cousinweb}

\begin{grammarlens}[title={What you've built}]
One more word for giving reasons and opinions.
\end{grammarlens}

\subsection*{Your turn}
\begin{itemize}
\raggedright
  \item \emph{Say it:} \emph{ghaḍvũ}
  \item \emph{Say it:} \emph{ghaḍvũ}, once more
  \item \emph{Say it:} \emph{sarkār navo kāydo ghaḍśe.} — The government will draw up a new law.
  \item \emph{Recall:} say the Gujarati for power to rule, then the Gujarati for the people of a state, then say \emph{ghaḍvũ} again
\end{itemize}

\begin{tcolorbox}[breakable,colback=teal!4,colframe=teal!35!black,title={Before you move on}]
What does \gu{ઘડવું} mean? (\textquotedblleft{}To shape; to frame (a law or a plan)\textquotedblright{}.) Say it once more.
\end{tcolorbox}
